En un ensayo idealizado, un rotor de momento de inercia \(J=2\,\mathrm{kg\,m^2}\) gira inicialmente a \(3\,\mathrm{rad/s}\). Al desconectar el motor, el único par que actúa es el de un freno viscoso: se opone al giro y vale \(-b\omega\), donde \(\omega\) es la velocidad angular y \(b=0.5\,\mathrm{N\,m\,s}\). El cambio del momento angular \(J\omega\) por unidad de tiempo es igual al par total. Se desprecia cualquier rozamiento adicional.

Tome como origen angular la posición del rotor al desconectar el motor y formule un modelo que permita saber qué ángulo recorre. Hay tres marcas situadas a \(8\), \(12\) y \(15\) radianes por delante de esa posición, contando vueltas completas. Determine cuáles alcanza y en qué instante cuando proceda. Dé después un criterio para una marca situada a un ángulo cualquiera \(\Theta>0\), y explique si la velocidad angular se anula en un tiempo finito.
